% !PRINTED 247
\weberpage{247}
If the functions $\Phi_{a,c,\partial}$ are so arranged, say by a
suitable determination of the constants that are still at our disposal
in them, that they remain finite for an infinite imaginary $\omega$ as
well, i.e. for $q=0$, then these functions are finite for an infinite
$j(\omega)$ as well, and the coefficients in (6) are constant. But
this is possible only if the $\Phi_{a,c,\partial}$ are \gesperrt{all equal to
one and the same constant $C$}. If one knows $\Phi_{a,c,\partial}$ as a
rational function of another quantity $\Psi_{a,c,\partial}$, then
$\Phi_{a,c,\partial} = C$ is either an \gesperrt{identity} or a \gesperrt{transformation
equation}
for $\Psi_{a,c,\partial}$.

\begin{center}
\textbf{§ 71. \gesperrt{The Transformation Equations for $\gamma_2$ and $\gamma_3$.}}
\end{center}

Transformation equations of the first level are obtained first of all
from the consideration of the functions § 54, (4), (5):
\[
\gamma_2(\omega) = \sqrt[3]{j(\omega)}, \quad \gamma_3(\omega) = \sqrt{j(\omega)-1728}.
\]
If $n$ is not divisible by 3, we can choose the numbers $a$, $c$,
$\partial$ in such a way that \gesperrt{$c$ always becomes divisible
by 3}.

Now, by § 54, (15), a linear transformation
$\begin{pmatrix} \alpha, & \beta \\ \gamma, & \delta \end{pmatrix}$
has the following effect on the function $\gamma_2(\omega)$:
\[
\gamma_2\left(\frac{\gamma+\delta\omega}{\alpha+\beta\omega}\right) = e^{-\frac{2\pi i}{3}(\alpha\gamma+\beta\delta-\alpha\beta-\alpha^2\beta\delta)} \gamma_2(\omega).
\]
In the compositions (8), (10) of § 69 we obtain
\[
\lambda \equiv n, \quad \delta \equiv 0 \ (\mathrm{mod}\ 3),
\]
and then $\alpha$ can be further determined in such a way that
$\alpha$ also, and consequently $c_2$, become divisible by 3.

From this it follows that by the two substitutions
\[
(\omega, \omega+1), \quad \left(\omega, -\frac{1}{\omega}\right)
\]
the $\nu$ functions
\weq{1}{\gamma_2\left(\frac{c+\partial\omega}{a}\right) \gamma_2(\omega)^{-n}}
are only permuted among one another, and hence are the roots of a
transformation equation of the first level.

If, secondly, $n$ is an \gesperrt{odd} number, let $c$ be assumed
\gesperrt{even}.
